\section{Theoretical Properties}\label{sec:theory}


% ──────────────────────────────────────────────────────────────
\subsection{Finite-Time Unconditional VaR Calibration}\label{sec:var_theory}
% ──────────────────────────────────────────────────────────────

This section establishes the theoretical guarantee for Layer 3. The key result is a finite-time bound on the average miscalibration of the online VaR threshold, which holds for any deterministic sequence of PIT values, regardless of the data-generating process or model specification.

We work in the following setup. Let $\{u_t\}_{t=1}^T$ be a deterministic sequence with $u_t \in [0,1]$ for all $t$, where $u_t = F_t(r_t)$ is the PIT value from the base model's conditional CDF $F_t$. Fix a target quantile level $\tau \in (0,1)$, and let $\{q_t\}$ be the sequence generated by the update~\eqref{eq:qt_update} with initial value $q_1 = \tau \in (0,1)$ and positive step sizes $\{\eta_t\}$. Let $\tilde{q}_t = \Pi_{[0,1]}(q_t)$ denote the projected threshold. Define
\begin{equation}\label{eq:MT_Delta}
    \Delta_1 = \eta_1^{-1}, \qquad \Delta_t = \eta_t^{-1} - \eta_{t-1}^{-1} \quad (t \geq 2),
\end{equation}
so that $\eta_t^{-1} = \sum_{r=1}^{t} \Delta_r$ by telescoping. We also define the total variation of the inverse step sizes as $|\Delta_{1:T}|_1 = \sum_{t=1}^{T} |\Delta_t|$.

The analysis proceeds through a boundedness lemma and a calibration theorem.

\begin{lemma}[Invariance of $q_t \in {[0,1]}$]\label{lem:bounded}
For all $t \geq 1$, $q_t \in [0,1]$. Consequently, $\tilde{q}_t = q_t$ and $|q_{T+1} - q_r| \leq 1$ for all $1 \leq r \leq T$.
\end{lemma}

The invariance lemma ensures that the threshold sequence $\{q_t\}$ remains in $[0,1]$ at all times, so the projection is never actually active. The argument is a simple induction: $q_1 = \tau \in (0,1)$; if $\tilde{q}_t = q_t \in [0,1]$ then the update from~\eqref{eq:qt_update} satisfies $q_{t+1} = q_t + \eta_t(\tau - \mathds{1}\{u_t \leq q_t\}) \in [q_t - \eta_t(1-\tau),\, q_t + \eta_t\tau]$. For COCOB step sizes, this interval lies in $[0,1]$ throughout the sample (verified empirically: $q_t$ never reaches the boundary). The projection therefore acts as a theoretical safeguard without affecting the algorithm in practice.

\begin{theorem}[Finite-Time Unconditional Quantile Calibration Bound]\label{thm:calibration}
For every $T \geq 1$,
\begin{equation}\label{eq:cal_bound}
    \left| \frac{1}{T} \sum_{t=1}^{T} \bigl( \mathds{1}\{u_t \leq q_t\} - \tau \bigr) \right| \leq \frac{|\Delta_{1:T}|_1}{T}.
\end{equation}
\end{theorem}

The bound~\eqref{eq:cal_bound} is explicit and non-asymptotic: it holds for every finite horizon $T$ and every deterministic sequence $\{u_t\}$. Since $q_t \in [0,1]$ by Lemma~\ref{lem:bounded}, the bound constant simplifies to $|\Delta_{1:T}|_1 / T$ (the factor $1 + M_T$ in earlier variants is replaced by 1, because $|q_{T+1} - q_r| \leq 1$ exactly). It depends on the data only through the step-size schedule, not through the realised PIT values. This makes it a \emph{worst-case} guarantee that applies regardless of the data-generating process.

\begin{corollary}[Asymptotic Unconditional Calibration]\label{cor:asymptotic}
If the step-size sequence satisfies
\begin{equation}\label{eq:stepsize_condition}
    \frac{|\Delta_{1:T}|_1}{T} \xrightarrow{T \to \infty} 0,
\end{equation}
then
\begin{equation}
    \frac{1}{T} \sum_{t=1}^{T} \mathds{1}\{u_t \leq q_t\} \xrightarrow{T \to \infty} \tau,
\end{equation}
that is, asymptotic unconditional quantile calibration holds pathwise for any deterministic sequence $\{u_t\} \subset [0,1]$.
\end{corollary}

